\documentclass[12pt,a4paper]{article}

% Packages
\usepackage[utf8]{inputenc}
\usepackage[T1]{fontenc}
\usepackage{geometry}
\usepackage{graphicx}
\usepackage{listings}
\usepackage{xcolor}
\usepackage{hyperref}
\usepackage{fancyhdr}
\usepackage{titlesec}
\usepackage{booktabs}
\usepackage{array}
\usepackage{longtable}
\usepackage{float}
\usepackage{caption}
\usepackage{subcaption}

% Page geometry
\geometry{
    left=2.5cm,
    right=2.5cm,
    top=2.5cm,
    bottom=2.5cm
}

% Colors for code listings
\definecolor{codegreen}{rgb}{0,0.6,0}
\definecolor{codegray}{rgb}{0.5,0.5,0.5}
\definecolor{codepurple}{rgb}{0.58,0,0.82}
\definecolor{backcolour}{rgb}{0.95,0.95,0.92}
\definecolor{jenkinsblue}{rgb}{0.2,0.4,0.6}

% Code listing style
\lstdefinestyle{mystyle}{
    backgroundcolor=\color{backcolour},
    commentstyle=\color{codegreen},
    keywordstyle=\color{magenta},
    numberstyle=\tiny\color{codegray},
    stringstyle=\color{codepurple},
    basicstyle=\ttfamily\footnotesize,
    breakatwhitespace=false,
    breaklines=true,
    captionpos=b,
    keepspaces=true,
    numbers=left,
    numbersep=5pt,
    showspaces=false,
    showstringspaces=false,
    showtabs=false,
    tabsize=2,
    frame=single,
    rulecolor=\color{black}
}
\lstset{style=mystyle}

% Groovy language definition for listings
\lstdefinelanguage{Groovy}{
    keywords={pipeline, agent, any, environment, stages, stage, steps, script, post, failure, fixed, sh, git, docker, mail, echo, def, if, else, return, true, false},
    keywordstyle=\color{blue}\bfseries,
    ndkeywords={String, Integer, Boolean},
    ndkeywordstyle=\color{darkgray}\bfseries,
    sensitive=true,
    comment=[l]{//},
    morecomment=[s]{/*}{*/},
    commentstyle=\color{codegreen}\ttfamily,
    stringstyle=\color{codepurple}\ttfamily,
    morestring=[b]',
    morestring=[b]"
}

% Header and footer
\pagestyle{fancy}
\fancyhf{}
\rhead{MT2025001}
\lhead{SPE Mini Project Report}
\cfoot{\thepage}

% Hyperref setup
\hypersetup{
    colorlinks=true,
    linkcolor=blue,
    filecolor=magenta,
    urlcolor=cyan,
}

% Title formatting
\titleformat{\section}
{\normalfont\Large\bfseries\color{jenkinsblue}}{\thesection}{1em}{}
\titleformat{\subsection}
{\normalfont\large\bfseries}{\thesubsection}{1em}{}

\begin{document}

% Title Page
\begin{titlepage}
    \centering
    \vspace*{2cm}

    {\Huge\bfseries Scientific Calculator with DevOps\\[0.5cm]}
    {\Large CI/CD Pipeline Implementation}

    \vspace{1.5cm}

    {\large\textbf{Mini Project Report}}\\
    {\large Software Production Engineering (CS816)}

    \vspace{2cm}

    \begin{figure}[H]
        \centering
        % Placeholder for institute logo
        % \includegraphics[width=0.3\textwidth]{iiitb_logo.png}
        \fbox{\parbox{4cm}{\centering\vspace{1cm}Institute Logo\\(Optional)\vspace{1cm}}}
    \end{figure}

    \vspace{1.5cm}

    {\Large\textbf{Author:}}\\[0.3cm]
    {\large Aayank Singhai}\\
    {\large Roll No: MT2025001}

    \vspace{1cm}

    {\Large\textbf{Submitted to:}}\\[0.3cm]
    {\large International Institute of Information Technology, Bangalore}

    \vfill

    {\large February 2026}

\end{titlepage}

% Table of Contents
\tableofcontents
\newpage

%==============================================================================
\section{Project Overview}
%==============================================================================

This project implements a \textbf{Scientific Calculator} application using Java Swing with a complete DevOps pipeline for automated build, test, and deployment. The project demonstrates the integration of modern CI/CD practices using Jenkins, Docker containerization, and Ansible-based deployment.

\subsection{Project Objectives}

\begin{itemize}
    \item Develop a functional scientific calculator with a modern GUI
    \item Implement comprehensive unit testing using JUnit
    \item Create an automated CI/CD pipeline using Jenkins
    \item Containerize the application using Docker with web-based access
    \item Automate deployment using Ansible
    \item Implement email notifications for build status
\end{itemize}

\subsection{Technology Stack}

\begin{table}[H]
    \centering
    \begin{tabular}{|l|l|l|}
        \hline
        \textbf{Category} & \textbf{Technology} & \textbf{Purpose} \\
        \hline
        Language & Java 11 & Application development \\
        GUI Framework & Java Swing & User interface \\
        Build Tool & Apache Maven & Dependency management \& build \\
        Testing & JUnit 4.13.2 & Unit testing \\
        CI/CD & Jenkins & Pipeline automation \\
        Containerization & Docker + NoVNC & Web-accessible container \\
        Deployment & Ansible & Automated deployment \\
        Version Control & Git/GitHub & Source code management \\
        \hline
    \end{tabular}
    \caption{Technology Stack}
\end{table}

%==============================================================================
\section{Application Architecture}
%==============================================================================

\subsection{Project Structure}

The project follows a standard Maven directory structure with clear separation between main source code and test code:

\begin{lstlisting}[language=bash, caption=Project Directory Structure]
CS816-SPE-MiniProject-Scientific_Calculator_With_DevOps/
|-- src/
|   |-- main/java/com/calculator/
|   |   |-- Calculator.java        # GUI Application
|   |   |-- CalculatorLogic.java   # Business Logic
|   |-- test/java/com/calculator/
|       |-- CalculatorLogicTest.java  # Unit Tests
|-- Dockerfile                     # Container configuration
|-- Jenkinsfile                    # CI/CD Pipeline
|-- deploy.yml                     # Ansible Playbook
|-- inventory                      # Ansible Inventory
|-- pom.xml                        # Maven Configuration
\end{lstlisting}

\subsection{Calculator Features}

The scientific calculator implements the following operations:

\begin{table}[H]
    \centering
    \begin{tabular}{|l|l|l|}
        \hline
        \textbf{Operation} & \textbf{Symbol} & \textbf{Constraints} \\
        \hline
        Addition & + & None \\
        Subtraction & - & Also handles negative numbers \\
        Multiplication & * & None \\
        Division & / & Divisor $\neq$ 0 \\
        Square Root & $\sqrt{}$ & Input $\geq$ 0 \\
        Factorial & n! & Non-negative integer $\leq$ 170 \\
        Natural Logarithm & ln & Input $>$ 0 \\
        Power & $x^y$ & Valid mathematical operation \\
        \hline
    \end{tabular}
    \caption{Calculator Operations}
\end{table}

%==============================================================================
\section{Jenkinsfile Explanation}
%==============================================================================

The CI/CD pipeline is defined in the \texttt{Jenkinsfile} using declarative pipeline syntax. Below is a detailed explanation of each component.

\subsection{Pipeline Configuration}

\begin{lstlisting}[language=Groovy, caption=Pipeline Environment Configuration]
pipeline {
    agent any
    environment {
        DOCKER_IMAGE_NAME = 'scientific-calculator'
        GITHUB_REPO_URL = 'https://github.com/aayanksinghai/
            CS816-SPE-MiniProject-Scientific_Calculator_With_DevOps'
        DOCKER_HUB_USERNAME = 'aayanksinghai'
        EMAIL_RECIPIENT = 'aayanksinghai02@gmail.com'
    }
\end{lstlisting}

\textbf{Explanation:}
\begin{itemize}
    \item \texttt{agent any}: Pipeline can run on any available Jenkins agent
    \item \texttt{environment}: Defines global environment variables used across all stages
    \item Variables store Docker image name, GitHub repository URL, Docker Hub credentials, and email recipient for notifications
\end{itemize}

\subsection{Stage 1: Clone Git Repository}

\begin{lstlisting}[language=Groovy, caption=Git Clone Stage]
stage('Clone Git') {
    steps {
        script {
            git branch: 'main',
                credentialsId: 'github_credentials',
                url: "${GITHUB_REPO_URL}"
        }
    }
}
\end{lstlisting}

\textbf{Purpose:} Fetches the latest source code from the GitHub repository.

\textbf{Key Elements:}
\begin{itemize}
    \item Clones the \texttt{main} branch
    \item Uses stored Jenkins credentials (\texttt{github\_credentials}) for authentication
    \item Repository URL is referenced from environment variables
\end{itemize}

\subsection{Stage 2: Build Maven Project}

\begin{lstlisting}[language=Groovy, caption=Maven Build Stage]
stage('Build the Maven Project') {
    steps {
        sh 'mvn clean package -DskipTests'
    }
}
\end{lstlisting}

\textbf{Purpose:} Compiles the Java source code and packages it into an executable JAR file.

\textbf{Maven Commands:}
\begin{itemize}
    \item \texttt{clean}: Removes previous build artifacts from \texttt{target/} directory
    \item \texttt{package}: Compiles source code and creates \texttt{scientific-calculator-1.0-SNAPSHOT.jar}
    \item \texttt{-DskipTests}: Skips tests during build (tests run in separate stage)
\end{itemize}

\subsection{Stage 3: Test Maven Project}

\begin{lstlisting}[language=Groovy, caption=Maven Test Stage]
stage('Test the Maven project') {
    steps {
        sh 'mvn test'
    }
}
\end{lstlisting}

\textbf{Purpose:} Executes all JUnit test cases to validate application functionality.

\textbf{Test Execution:}
\begin{itemize}
    \item Runs all tests in \texttt{src/test/java/}
    \item Generates test reports in \texttt{target/surefire-reports/}
    \item Pipeline fails if any test fails
\end{itemize}

\subsection{Stage 4: Build Docker Image}

\begin{lstlisting}[language=Groovy, caption=Docker Build Stage]
stage('Build Docker Image') {
    steps {
        script {
            docker.build("${DOCKER_IMAGE_NAME}", '.')
        }
    }
}
\end{lstlisting}

\textbf{Purpose:} Creates a Docker image containing the calculator application with NoVNC for web access.

\textbf{Docker Build Process:}
\begin{itemize}
    \item Uses \texttt{Dockerfile} in project root
    \item Image includes: OpenJDK 17, Xvfb, x11vnc, NoVNC, Openbox
    \item Enables GUI application to run in containerized environment
\end{itemize}

\subsection{Stage 5: Push to Docker Hub}

\begin{lstlisting}[language=Groovy, caption=Docker Push Stage]
stage('Push Docker Image to Docker Hub') {
    steps {
        script {
            docker.withRegistry('', 'DockerHubCred') {
                sh "docker tag ${DOCKER_IMAGE_NAME}
                    ${DOCKER_HUB_USERNAME}/${DOCKER_IMAGE_NAME}:latest"
                sh "docker push
                    ${DOCKER_HUB_USERNAME}/${DOCKER_IMAGE_NAME}:latest"
            }
        }
    }
}
\end{lstlisting}

\textbf{Purpose:} Pushes the built Docker image to Docker Hub registry for distribution.

\textbf{Process:}
\begin{itemize}
    \item Authenticates with Docker Hub using stored credentials
    \item Tags image with Docker Hub username and \texttt{latest} tag
    \item Pushes image to \texttt{aayanksinghai/scientific-calculator:latest}
\end{itemize}

\subsection{Stage 6: Deploy with Ansible}

\begin{lstlisting}[language=Groovy, caption=Ansible Deployment Stage]
stage('Deploy with Ansible') {
    steps {
        sh 'ansible-playbook -i inventory deploy.yml'
    }
}
\end{lstlisting}

\textbf{Purpose:} Automates application deployment using Ansible playbook.

\textbf{Ansible Playbook Actions:}
\begin{enumerate}
    \item Pulls latest Docker image from Docker Hub
    \item Stops and removes existing container (if any)
    \item Starts new container with port mapping (8081:8080)
    \item Verifies container status
\end{enumerate}

\subsection{Post-Build Actions: Email Notifications}

\begin{lstlisting}[language=Groovy, caption=Email Notification Configuration]
post {
    failure {
        script {
            mail to: "${EMAIL_RECIPIENT}",
                 subject: "BUILD FAILED: ${env.JOB_NAME} -
                          Build #${env.BUILD_NUMBER}",
                 body: """Build Failed!
                 Project: ${env.JOB_NAME}
                 Build Number: ${env.BUILD_NUMBER}
                 Build URL: ${env.BUILD_URL}
                 Please check the console output."""
        }
    }
    fixed {
        script {
            mail to: "${EMAIL_RECIPIENT}",
                 subject: "BUILD FIXED: ${env.JOB_NAME} -
                          Build #${env.BUILD_NUMBER}",
                 body: """Build Fixed!
                 The build is back to normal."""
        }
    }
}
\end{lstlisting}

\textbf{Purpose:} Sends email notifications based on build status.

\textbf{Notification Triggers:}
\begin{itemize}
    \item \texttt{failure}: Sends email when build fails at any stage
    \item \texttt{fixed}: Sends email when a previously failed build succeeds
\end{itemize}

%==============================================================================
\section{Pipeline Flow Diagram}
%==============================================================================

\begin{figure}[H]
    \centering
    \fbox{\parbox{0.9\textwidth}{\centering\vspace{2cm}
    \textbf{[SCREENSHOT PLACEHOLDER]}\\[0.5cm]
    Insert Jenkins Pipeline Flow Diagram or\\
    Blue Ocean Pipeline Visualization here\\[0.5cm]
    Suggested: Screenshot of Jenkins Blue Ocean showing all 6 stages
    \vspace{2cm}}}
    \caption{Jenkins Pipeline Stages Visualization}
    \label{fig:pipeline-flow}
\end{figure}

The pipeline executes in the following sequential order:

\begin{center}
\fbox{Clone Git} $\rightarrow$ \fbox{Maven Build} $\rightarrow$ \fbox{Maven Test} $\rightarrow$ \fbox{Docker Build} $\rightarrow$ \fbox{Push to Hub} $\rightarrow$ \fbox{Ansible Deploy}
\end{center}

%==============================================================================
\section{Test Results}
%==============================================================================

\subsection{Test Cases Overview}

The test suite \texttt{CalculatorLogicTest.java} contains comprehensive unit tests for all calculator operations:

\begin{table}[H]
    \centering
    \begin{tabular}{|l|l|p{6cm}|}
        \hline
        \textbf{Test Method} & \textbf{Status} & \textbf{Description} \\
        \hline
        testAddition & \textcolor{green}{PASS} & Tests addition with positive, negative, and decimal numbers \\
        testSubtraction & \textcolor{green}{PASS} & Tests subtraction operations \\
        testMultiplication & \textcolor{green}{PASS} & Tests multiplication including zero \\
        testDivision & \textcolor{green}{PASS} & Tests division with various inputs \\
        testDivisionByZero & \textcolor{green}{PASS} & Verifies ArithmeticException is thrown \\
        testZeroOperations & \textcolor{green}{PASS} & Tests operations involving zero \\
        testSquareRoot & \textcolor{green}{PASS} & Tests square root of positive numbers \\
        testSquareRootNegative & \textcolor{green}{PASS} & Verifies exception for negative input \\
        testFactorial & \textcolor{green}{PASS} & Tests factorial (0! to 10!) \\
        testFactorialNegative & \textcolor{green}{PASS} & Verifies exception for negative input \\
        testFactorialNonInteger & \textcolor{green}{PASS} & Verifies exception for decimal input \\
        testNaturalLog & \textcolor{green}{PASS} & Tests ln(1), ln(e), ln(10), ln(2) \\
        testNaturalLogZero & \textcolor{green}{PASS} & Verifies exception for zero \\
        testNaturalLogNegative & \textcolor{green}{PASS} & Verifies exception for negative input \\
        testPower & \textcolor{green}{PASS} & Tests power function with various exponents \\
        \hline
    \end{tabular}
    \caption{Unit Test Results Summary}
\end{table}

\subsection{Test Execution Screenshot}

\begin{figure}[H]
    \centering
    \fbox{\parbox{0.9\textwidth}{\centering\vspace{3cm}
    \textbf{[SCREENSHOT PLACEHOLDER]}\\[0.5cm]
    Insert screenshot of Maven test execution output showing:\\
    - All 15 tests passing\\
    - Test execution time\\
    - BUILD SUCCESS message
    \vspace{3cm}}}
    \caption{Maven Test Execution Results}
    \label{fig:test-results}
\end{figure}

\subsection{Test Observations}

\begin{itemize}
    \item \textbf{Total Tests:} 15 test methods covering all calculator operations
    \item \textbf{Pass Rate:} 100\% (all tests pass successfully)
    \item \textbf{Exception Handling:} 6 tests verify proper exception handling for invalid inputs
    \item \textbf{Edge Cases:} Tests include boundary conditions like factorial of 0, ln(1), and division by zero
    \item \textbf{Precision:} Floating-point comparisons use delta of 0.001 for accuracy
\end{itemize}

%==============================================================================
\section{Pipeline Execution Screenshots}
%==============================================================================

\subsection{Successful Pipeline Run}

\begin{figure}[H]
    \centering
    \fbox{\parbox{0.9\textwidth}{\centering\vspace{3cm}
    \textbf{[SCREENSHOT PLACEHOLDER]}\\[0.5cm]
    Insert screenshot of successful Jenkins pipeline run showing:\\
    - All 6 stages completed (green checkmarks)\\
    - Build duration\\
    - Console output summary
    \vspace{3cm}}}
    \caption{Successful Pipeline Execution}
    \label{fig:success-pipeline}
\end{figure}

\subsection{Docker Hub Image}

\begin{figure}[H]
    \centering
    \fbox{\parbox{0.9\textwidth}{\centering\vspace{2.5cm}
    \textbf{[SCREENSHOT PLACEHOLDER]}\\[0.5cm]
    Insert screenshot of Docker Hub showing:\\
    - Repository: aayanksinghai/scientific-calculator\\
    - Latest tag with push timestamp\\
    - Image size information
    \vspace{2.5cm}}}
    \caption{Docker Hub Repository}
    \label{fig:dockerhub}
\end{figure}

\subsection{Application Running in Container}

\begin{figure}[H]
    \centering
    \fbox{\parbox{0.9\textwidth}{\centering\vspace{3cm}
    \textbf{[SCREENSHOT PLACEHOLDER]}\\[0.5cm]
    Insert screenshot of calculator application accessed via:\\
    http://localhost:8081/vnc.html\\[0.3cm]
    Show the calculator GUI with dark theme and buttons
    \vspace{3cm}}}
    \caption{Calculator Application via NoVNC}
    \label{fig:calculator-app}
\end{figure}

\subsection{Calculator UI Demo}

The calculator features a modern dark-themed GUI with color-coded buttons for enhanced user experience:

\begin{figure}[H]
    \centering
    \fbox{\parbox{0.9\textwidth}{\centering\vspace{2.5cm}
    \textbf{[SCREENSHOT PLACEHOLDER]}\\[0.5cm]
    Insert screenshot of Calculator UI showing:\\
    - Dark theme interface\\
    - Color-coded buttons (Orange: operators, Blue: scientific, Red: clear)\\
    - Expression display and result display
    \vspace{2.5cm}}}
    \caption{Calculator User Interface}
    \label{fig:calculator-ui}
\end{figure}

\subsubsection{UI Features Demonstrated}

\begin{itemize}
    \item \textbf{Expression Display:} Secondary display showing current expression (e.g., "5 + 3 =")
    \item \textbf{Result Display:} Primary display showing input numbers and calculation results
    \item \textbf{Color Scheme:}
    \begin{itemize}
        \item Orange buttons: Arithmetic operators (+, -, *, /, =)
        \item Blue buttons: Scientific functions ($\sqrt{}$, n!, ln, $x^y$)
        \item Red button: Clear (C)
        \item Gray buttons: Number pad (0-9) and decimal point
    \end{itemize}
    \item \textbf{Error Handling:} Red text display for error messages (e.g., "Can't divide by 0")
\end{itemize}

\subsubsection{Sample Calculations Demo}

\begin{figure}[H]
    \centering
    \fbox{\parbox{0.9\textwidth}{\centering\vspace{2.5cm}
    \textbf{[SCREENSHOT PLACEHOLDER]}\\[0.5cm]
    Insert screenshot(s) showing sample calculations:\\
    - Basic operation: 25 + 17 = 42\\
    - Scientific operation: $\sqrt{144}$ = 12\\
    - Factorial: 5! = 120\\
    - Error handling: Division by zero error message
    \vspace{2.5cm}}}
    \caption{Calculator Operations Demo}
    \label{fig:calculator-demo}
\end{figure}

\subsection{Email Notification (Build Failure)}

\begin{figure}[H]
    \centering
    \fbox{\parbox{0.9\textwidth}{\centering\vspace{2.5cm}
    \textbf{[SCREENSHOT PLACEHOLDER]}\\[0.5cm]
    Insert screenshot of email received on build failure showing:\\
    - Subject: BUILD FAILED: scientific-calculator - Build \#N\\
    - Email body with project details and build URL
    \vspace{2.5cm}}}
    \caption{Build Failure Email Notification}
    \label{fig:email-failure}
\end{figure}

%==============================================================================
\section{GitHub Webhook Integration}
%==============================================================================

The project integrates GitHub Webhooks with Jenkins to enable automatic build triggers whenever code is pushed to the repository. This implements a true Continuous Integration workflow where every commit automatically initiates the pipeline.

\subsection{Webhook Configuration}

The webhook is configured in the GitHub repository settings to notify Jenkins on push events:

\begin{itemize}
    \item \textbf{Payload URL:} \texttt{http://<jenkins-server>:8080/github-webhook/}
    \item \textbf{Content Type:} \texttt{application/json}
    \item \textbf{Trigger Event:} Push events to the repository
    \item \textbf{Branch Filter:} \texttt{main} branch
\end{itemize}

\subsection{Jenkins SCM Configuration}

In the Jenkins job configuration, GitHub hook trigger is enabled:

\begin{itemize}
    \item \textbf{Build Trigger:} GitHub hook trigger for GITScm polling
    \item \textbf{SCM:} Git with repository URL and credentials
    \item \textbf{Branch Specifier:} \texttt{*/main}
\end{itemize}

\subsection{Webhook Flow}

\begin{enumerate}
    \item Developer pushes code to GitHub repository
    \item GitHub sends POST request to Jenkins webhook endpoint
    \item Jenkins receives the webhook payload and validates the request
    \item Jenkins triggers the pipeline automatically
    \item Pipeline executes all stages (Clone, Build, Test, Docker, Deploy)
    \item Build status is reported back and email notifications are sent if needed
\end{enumerate}

\subsection{Webhook Configuration Screenshot}

\begin{figure}[H]
    \centering
    \fbox{\parbox{0.9\textwidth}{\centering\vspace{2.5cm}
    \textbf{[SCREENSHOT PLACEHOLDER]}\\[0.5cm]
    Insert screenshot of GitHub Webhook configuration showing:\\
    - Payload URL pointing to Jenkins\\
    - Content type: application/json\\
    - Recent deliveries with green checkmarks
    \vspace{2.5cm}}}
    \caption{GitHub Webhook Configuration}
    \label{fig:webhook-config}
\end{figure}

\subsection{Automatic Build Trigger Screenshot}

\begin{figure}[H]
    \centering
    \fbox{\parbox{0.9\textwidth}{\centering\vspace{2.5cm}
    \textbf{[SCREENSHOT PLACEHOLDER]}\\[0.5cm]
    Insert screenshot of Jenkins build history showing:\\
    - Build triggered by: "Started by GitHub push by aayanksinghai"\\
    - Timestamp of automatic trigger
    \vspace{2.5cm}}}
    \caption{Jenkins Build Triggered by GitHub Webhook}
    \label{fig:webhook-trigger}
\end{figure}

%==============================================================================
\section{Ansible Deployment}
%==============================================================================

The Ansible playbook (\texttt{deploy.yml}) automates the deployment process:

\begin{lstlisting}[language=bash, caption=Ansible Playbook (deploy.yml)]
---
- name: Deploy Calculator App with Docker
  hosts: webservers
  tasks:
    - name: Pull the latest Docker image
      shell: docker pull aayanksinghai/scientific-calculator:latest

    - name: Stop and remove existing container
      shell: docker rm -f scientific-calculator
      ignore_errors: true

    - name: Run calculator container (detached)
      shell: docker run -d -p 8081:8080 --name scientific-calculator
             --restart unless-stopped
             aayanksinghai/scientific-calculator:latest

    - name: Display container status
      shell: docker ps -f "name=scientific-calculator"
\end{lstlisting}

\textbf{Deployment Process:}
\begin{enumerate}
    \item Pulls the latest image from Docker Hub
    \item Removes any existing container to avoid conflicts
    \item Starts a new container with:
    \begin{itemize}
        \item Port mapping: Host 8081 $\rightarrow$ Container 8080
        \item Auto-restart policy: \texttt{unless-stopped}
        \item Detached mode for background execution
    \end{itemize}
    \item Verifies deployment by checking container status
\end{enumerate}

%==============================================================================
\section{Conclusion}
%==============================================================================

This project successfully demonstrates the implementation of a complete DevOps pipeline for a Java-based scientific calculator application. Key achievements include:

\begin{itemize}
    \item \textbf{Automated CI/CD Pipeline:} 6-stage Jenkins pipeline covering build, test, containerization, and deployment
    \item \textbf{Comprehensive Testing:} 15 unit tests with 100\% pass rate covering all calculator operations and edge cases
    \item \textbf{Docker Containerization:} Web-accessible GUI application using NoVNC technology
    \item \textbf{Automated Deployment:} Ansible-based deployment ensuring consistent and repeatable deployments
    \item \textbf{Monitoring:} Email notifications for build failures and recovery
\end{itemize}

The project showcases practical application of DevOps principles including:
\begin{itemize}
    \item Continuous Integration (automated builds and tests)
    \item Continuous Deployment (automated container deployment)
    \item Infrastructure as Code (Jenkinsfile, Dockerfile, Ansible playbook)
    \item Version Control (Git-based workflow)
\end{itemize}

%==============================================================================
\section*{References}
%==============================================================================

\begin{itemize}
    \item Jenkins Documentation: \url{https://www.jenkins.io/doc/}
    \item Docker Documentation: \url{https://docs.docker.com/}
    \item Ansible Documentation: \url{https://docs.ansible.com/}
    \item Maven Documentation: \url{https://maven.apache.org/guides/}
    \item JUnit 4 Documentation: \url{https://junit.org/junit4/}
    \item NoVNC Project: \url{https://novnc.com/}
\end{itemize}

\vspace{1cm}
\hrule
\vspace{0.5cm}
\begin{center}
    \textbf{GitHub Repository:} \url{https://github.com/aayanksinghai/CS816-SPE-MiniProject-Scientific_Calculator_With_DevOps}\\[0.3cm]
    \textbf{Docker Hub:} \url{https://hub.docker.com/r/aayanksinghai/scientific-calculator}
\end{center}

\end{document}

